\newpage

\cleardoublepage
\phantomsection
\addcontentsline{toc}{chapter}{参考文献}
\markboth{参考文献}{参考文献}

\begingroup
\sloppy
\setlength{\emergencystretch}{3em}

\begin{thebibliography}{999}

\bibitem{zhouyi} \textbf{《周易》}（含《易经》与《易传》）。通行本。

封面引用的“无平不陂，无往不复”出自《周易·泰卦》九三爻辞。阴阳消长、往复循环的思想是周期数学“对称性 $\to$ 周期性”纲领的古老回响。

\bibitem{daodejing} \textbf{《道德经》}，老子 著。通行本（王弼注本）。

封面引用的“道生一，一生二，二生三，三生万物”出自第四十二章。

\bibitem{zhuangzi} \textbf{《庄子》}，庄周 著。通行本。

“方生方死，方死方生”（《齐物论》）揭示了对立面的相互转化。

\bibitem{huangdineijing} \textbf{《黄帝内经》}（含《素问》与《灵枢》）。通行本。

五运六气体系是周期嵌套结构的早期模型。

\bibitem{shangshu} \textbf{《尚书》}（含《尧典》）。通行本。

“期三百有六旬有六日，以闰月定四时成岁”，是最早的阴阳合历官方记载。

\bibitem{huainanzi} \textbf{《淮南子》}，刘安 编。通行本。

序言引《天文训》及《览冥训》，作为轩辕历法建元和女娲补天历法解读的核心文本。

\bibitem{shanhaijing} \textbf{《山海经》}，佚名。通行本。

序言引精卫填海及夸父逐日，作为历法修复永恒性与观测先河的神话文本。

\bibitem{shiji} \textbf{《史记》}，司马迁 著。

《天官书》为古代天文观测的系统记录。《历书》载太初元年甲寅为轩辕历法建元推演的基准。

\bibitem{shuowen} \textbf{《说文解字》}，许慎 著。

序言引“闰，余分之月”，为后羿射日所留“一日”作为闰余的训诂学依据。

\bibitem{zhoubi} \textbf{《周髀算经》}，佚名。

“七衡六间”图用同心圆描述太阳周年视运动。

\bibitem{yueling} \textbf{《礼记·月令》}，佚名。通行本。

按月记载天象、物候与农事的周期节律，是最早的系统性物候学文献。

\bibitem{mengzi} \textbf{《孟子》}，孟轲 著。通行本。

“万物皆备于我矣，反身而诚，乐莫大焉”（《尽心上》）。

\bibitem{sunzi} \textbf{《孙子兵法》}，孙武 著。通行本。

“善战者，求之于势，不责于人”（《兵势》）。

\bibitem{shennong} \textbf{《神农本草经》}，佚名。通行本（孙星衍、孙冯翼辑本）。

最早的中药学经典，记载了365种药物的采收时节和产地。

\bibitem{bencao} \textbf{《本草纲目》}，李时珍 著（1596）。

载药1892种，系统记录了每味药物的“时月”、“土地”和“炮制”。

\bibitem{huangjijingshi} \textbf{《皇极经世》}，邵雍 著。

元会运世的嵌套周期体系是本书历法章周期嵌套结构的理论基础。

\bibitem{leibniz1703} \textbf{Gottfried Wilhelm Leibniz}. \textit{Explication de l'Arithm\'etique Binaire} (1703).

莱布尼茨二进制算术的原始论文。二进制列周期性（Periodus numerorum）的发现为本书位宽与带宽概念提供了历史前驱。

\bibitem{leibniz1679} \textbf{Gottfried Wilhelm Leibniz}. \textit{De Progressione Dyadica} (1679).

莱布尼茨二进制级数展开的手稿。自然数序列每一位的周期变化规律在此手稿中首次被系统记录。

\bibitem{leibniz1714} \textbf{Gottfried Wilhelm Leibniz}. \textit{Monadologie} (1714).

单子论。二进制从“无”（$0$）中生“有”（$1$）的思想与单子论中“从简单实体构造复杂世界”的哲学一脉相承。

\bibitem{galileo1623} \textbf{Galileo Galilei}. \textit{The Assayer} (1623).

封面引用的“数学是上帝用来书写宇宙的语言”出自此著。

\bibitem{newton1687} \textbf{Isaac Newton}. \textit{Philosophiae Naturalis Principia Mathematica} (1687).

万有引力定律和运动三定律为经典周期运动提供了第一个完整的数学描述。

\bibitem{euler1748} \textbf{Leonhard Euler}. \textit{Introductio in Analysin Infinitorum} (1748).

欧拉公式 $e^{i\theta} = \cos\theta + i\sin\theta$ 是周转圆模型和傅里叶分析的数学基石。

\bibitem{fourier1822} \textbf{Joseph Fourier}. \textit{Th\'eorie Analytique de la Chaleur} (1822).

傅里叶级数与傅里叶变换的原始著作。

\bibitem{hamilton1853} \textbf{William R. Hamilton}. \textit{Lectures on Quaternions}. Hodges and Smith, 1853.

四元数的原创性著作。

\bibitem{riemann1854} \textbf{Bernhard Riemann}. \textit{\"Uber die Hypothesen, welche der Geometrie zu Grunde liegen} (1854).

黎曼几何的奠基性演讲。

\bibitem{maxwell1873} \textbf{James Clerk Maxwell}. \textit{A Treatise on Electricity and Magnetism} (1873).

麦克斯韦方程组是本书电磁系统之周期结构章电磁场理论的物理原型。

\bibitem{klein1872} \textbf{Felix Klein}. \textit{Vergleichende Betrachtungen \"uber neuere geometrische Forschungen} (Erlangen Programme) (1872).

Klein的Erlangen纲领——几何学是研究变换群不变量的科学。

\bibitem{poincare1892} \textbf{Henri Poincar\'e}. \textit{Les M\'ethodes Nouvelles de la M\'ecanique C\'eleste} (1892--1899).

庞加莱对三体问题周期轨道的开创性研究。

\bibitem{planck1901} \textbf{Max Planck}. \textit{\"Uber das Gesetz der Energieverteilung im Normalspectrum}. \textit{Annalen der Physik}, 309(3), 553--563 (1901).

黑体辐射的普朗克公式原始论文。

\bibitem{noether1918} \textbf{Emmy Noether}. \textit{Invariante Variationsprobleme}. \textit{Nachrichten von der Gesellschaft der Wissenschaften zu G\"ottingen}, 235--257 (1918).

诺特定理的原始论文。每一个连续对称性对应一个守恒量——本书周期基础章守恒律与对称性关系的理论来源。

\bibitem{einstein1916} \textbf{Albert Einstein}. \textit{Die Grundlage der allgemeinen Relativit\"atstheorie}. \textit{Annalen der Physik}, 49, 769--822 (1916).

局域对称性决定引力动力学的思想。

\bibitem{compton1923} \textbf{Arthur H. Compton}. \textit{A Quantum Theory of the Scattering of X-rays by Light Elements}. \textit{Physical Review}, 21(5), 483--502 (1923).

康普顿散射的原始论文。

\bibitem{yangmills1954} \textbf{Chen Ning Yang, Robert L. Mills}. \textit{Conservation of Isotopic Spin and Isotopic Gauge Invariance}. \textit{Physical Review}, 96(1), 191--195 (1954).

杨-米尔斯规范场论的开创性论文。

\bibitem{weyl1928} \textbf{Hermann Weyl}. \textit{The Theory of Groups and Quantum Mechanics} (1928).

群表示论在量子力学中的首次系统应用。

\bibitem{weyl1950} \textbf{Hermann Weyl}. \textit{The Theory of Groups and Quantum Mechanics}. Dover Publications, 1950.

Weyl 群表示论经典著作的英译本。书中对紧致群表示论、特征标理论与对称群的系统处理，为本书群与群函数章、高阶群章提供了表示论的经典框架。

\bibitem{serre1977} \textbf{Jean-Pierre Serre}. \textit{Linear Representations of Finite Groups}. Springer, 1977.

群表示论的经典教材。有限群的特征标理论与诱导表示理论在本书周期函数与具有周期轨道的群章中被系统使用。

\bibitem{mackey1968} \textbf{George W. Mackey}. \textit{Induced Representations of Groups and Quantum Mechanics}. Benjamin, 1968.

诱导表示理论。Mackey 的诱导表示理论为本书周期范畴章中群作用的函子化表述提供了代数基础。

\bibitem{wigner1959} \textbf{Eugene P. Wigner}. \textit{Group Theory and Its Application to the Quantum Mechanics of Atomic Spectra}. Academic Press, 1959.

Wigner D-矩阵是本书SU(2)上调和分析的核心工具。

\bibitem{vonneumann1955} \textbf{John von Neumann}. \textit{Mathematical Foundations of Quantum Mechanics}. Princeton University Press, 1955.

量子力学的希尔伯特空间表述。书中对自伴算子谱定理与酉群的系统处理，为本书线性代数章谱定理与周期范畴章提供了算子理论的经典基础。

\bibitem{dirac1958} \textbf{Paul A. M. Dirac}. \textit{The Principles of Quantum Mechanics}. 4th ed., Oxford University Press, 1958.

狄拉克的变换理论。

\bibitem{weinstein1978} \textbf{Alan Weinstein}. \textit{Periodic orbits for convex Hamiltonian systems}. Annals of Mathematics, 108(3):507--518, 1978.
\bibitem{arnold1989} \textbf{Vladimir I. Arnold}. \textit{Mathematical Methods of Classical Mechanics}. 2nd ed., Springer, 1989.

作用量原理、辛几何、可积系统与KAM理论。Arnold 对刘维尔–阿诺德定理与不变环面的系统处理是本书周期流形章的核心参考。

\bibitem{arnold1988} \textbf{V. I. Arnold, V. V. Kozlov, A. I. Neishtadt}. \textit{Dynamical Systems III: Mathematical Aspects of Classical and Celestial Mechanics}. Springer, 1988.

KAM 理论与周期轨道扰动理论的系统阐述，本书周期基础章与天文系统章的重要参考。

\bibitem{meyer1992} \textbf{Kenneth R. Meyer, Glen R. Hall}. \textit{Introduction to Hamiltonian Dynamical Systems and the N-Body Problem}. Springer, 1992.

哈密顿动力系统与周期轨道续接、KAM 理论的经典教材。

\bibitem{perko1991} \textbf{Lawrence Perko}. \textit{Differential Equations and Dynamical Systems}. Springer, 1991.

本书全局理论章对周期轨道、极限环、Poincaré 映射与 Poincaré–Bendixson 定理的系统论述，为本书周期基础章闭合轨道判据与生物系统章极限环分析提供了标准参考。

\bibitem{iooss1990} \textbf{G\'erard Iooss, Daniel D. Joseph}. \textit{Elementary Stability and Bifurcation Theory}. 2nd ed., Springer, 1990.

Hopf 分岔与周期解分支的经典教材。书中对二维与一般情形下从稳态到周期解的分支、次谐波分支与准周期解的讨论，为本书周期基础章周期解生成机制与生物系统章振荡涌现提供了分岔论基础。

\bibitem{chicone2006} \textbf{Carmen Chicone}. \textit{Ordinary Differential Equations with Applications}. 2nd ed., Springer, 2006.

常微分方程及其应用，含 Floquet 理论与周期解的系统处理。

\bibitem{hirsch2013} \textbf{Morris W. Hirsch, Stephen Smale, Robert L. Devaney}. \textit{Differential Equations, Dynamical Systems, and an Introduction to Chaos}. 3rd ed., Academic Press, 2013.

微分方程、动力系统与混沌的经典教材，含周期轨道与极限集。

\bibitem{ruelle1989} \textbf{David Ruelle}. \textit{Elements of Differentiable Dynamics and Bifurcation Theory}. Academic Press, 1989.

可微动力系统、不动点与周期轨道、稳定与不稳定流形。

\bibitem{robinson2004} \textbf{R. Clark Robinson}. \textit{An Introduction to Dynamical Systems: Continuous and Discrete}. Pearson Prentice Hall, 2004.

连续与离散动力系统的入门教材，含周期轨道与混沌吸引子。

\bibitem{strogatz2015} \textbf{Steven H. Strogatz}. \textit{Nonlinear Dynamics and Chaos: With Applications to Physics, Biology, Chemistry, and Engineering}. 2nd ed., Westview Press, 2015.

非线性动力学与混沌的经典教材，极限环与周期振荡的标准参考。

\bibitem{kuchment1993} \textbf{Peter Kuchment}. \textit{Floquet Theory for Partial Differential Equations}. Birkh\"auser, 1993.

偏微分方程 Floquet 理论的权威著作。书中对周期系数线性微分方程的 Floquet–Lyapunov 定理、周期微分算子的谱理论以及周期偏微分方程在量子力学、流体力学与弹性理论中的应用的系统论述，构成本书微分方程与复变函数章 Floquet 理论的核心参考。

\bibitem{brown2013} \textbf{B. Malcolm Brown, Michael S. P. Eastham, Karl Michael Schmidt}. \textit{Periodic Differential Operators}. Springer, 2013.

周期微分算子的 Floquet 理论、振荡与谱分析。书中对一维周期系统从 Floquet 理论到谱分析的完整处理，为本书周期系数微分方程与周期流形章的 Bloch 波理论提供了标准参考。

\bibitem{steinweiss1971} \textbf{Elias M. Stein, Guido Weiss}. \textit{Introduction to Fourier Analysis on Euclidean Spaces}. Princeton University Press, 1971.

欧氏空间上的调和分析标准参考。

\bibitem{folland1995} \textbf{Gerald B. Folland}. \textit{A Course in Abstract Harmonic Analysis}. CRC Press, 1995.

局部紧群上的调和分析理论。Folland 对 Pontryagin 对偶、Peter–Weyl 定理与群上傅里叶变换的系统处理，为本书调和分析章的统一框架提供了标准参考。

\bibitem{rudin1962} \textbf{Walter Rudin}. \textit{Fourier Analysis on Groups}. Interscience, 1962.

局部紧阿贝尔群上调和分析的经典著作。

\bibitem{katznelson2004} \textbf{Yitzhak Katznelson}. \textit{An Introduction to Harmonic Analysis}. 3rd ed., Cambridge University Press, 2004.

调和分析的入门教材。含周期函数的傅里叶级数、$L^2$ 收敛与一致收敛。

\bibitem{zygmund1959} \textbf{Antoni Zygmund}. \textit{Trigonometric Series}. 2nd ed., Cambridge University Press, 1959.

三角级数理论的经典巨著。

\bibitem{grafakos2009} \textbf{Loukas Grafakos}. \textit{Modern Fourier Analysis}. 2nd ed., Springer, 2009.

现代傅里叶分析的系统教材，含周期函数多变量三角展开的历史渊源。

\bibitem{varadhan2000} \textbf{S. R. S. Varadhan}. \textit{Harmonic Analysis}. Courant Institute, 2000.

调和分析入门讲义，以周期函数分解为正弦余弦组合为核心范例。

\bibitem{pontryagin1939} \textbf{Lev S. Pontryagin}. \textit{Topological Groups}. Princeton University Press, 1939.

Pontryagin对偶定理的原始著作。本书调和分析章Pontryagin对偶定理的理论基础。

\bibitem{peterweyl1927} \textbf{Fritz Peter, Hermann Weyl}. \textit{Die Vollst\"andigkeit der primitiven Darstellungen einer geschlossenen kontinuierlichen Gruppe}. \textit{Mathematische Annalen}, 97(1), 737--755 (1927).

Peter-Weyl定理的原始论文。

\bibitem{fejer1904} \textbf{Leopold Fej\'er}. \textit{Untersuchungen \"uber Fouriersche Reihen}. \textit{Mathematische Annalen}, 58, 51--69 (1904).

Fej\'er核的原始论文。

\bibitem{bohr1925} \textbf{Harald Bohr}. \textit{Zur Theorie der fastperiodischen Funktionen}. \textit{Acta Mathematica}, 45, 29--127 (1925).

几乎周期函数的开创性论文，本书周期函数章几乎周期概念的理论来源。

\bibitem{corduneanu1989} \textbf{C. Corduneanu}. \textit{Almost Periodic Functions}. 2nd English ed., Chelsea, 1989.

几乎周期函数的系统教材，含玻尔均值与具有周期轨道的群结构的经典处理。

\bibitem{terras1999} \textbf{Audrey Terras}. \textit{Fourier Analysis on Finite Groups and Applications}. Cambridge University Press, 1999.

有限群傅里叶分析的入门教材，为本书循环群 $C_T$ 上的相位算子与傅里叶矩阵提供群表示论基础。

\bibitem{luong2009} \textbf{Bao Luong}. \textit{Fourier Analysis on Finite Abelian Groups}. Birkh\"auser, 2009.

有限阿贝尔群上的傅里叶分析，含周期函数的傅里叶变换与反演公式。

\bibitem{shannon1948} \textbf{Claude E. Shannon}. \textit{A Mathematical Theory of Communication}. \textit{Bell System Technical Journal}, 27, 379--423, 623--656 (1948).

信息熵与采样定理的奠基性著作。

\bibitem{shannon1949} \textbf{Claude E. Shannon}. \textit{Communication in the Presence of Noise}. \textit{Proceedings of the IRE}, 37(1), 10--21 (1949).

奈奎斯特-香农采样定理的严格证明。本书周期函数章香农-奈奎斯特采样定理的理论来源。

\bibitem{slepian1978} \textbf{David Slepian}. \textit{Prolate Spheroidal Wave Functions, Fourier Analysis, and Uncertainty---V: The Discrete Case}. \textit{Bell System Technical Journal}, 57(5), 1371--1430 (1978).

离散扁球（Slepian）序列的经典论文，本书多窗谱估计与能量集中性讨论的理论基础。

\bibitem{percival1993} \textbf{Donald B. Percival, Andrew T. Walden}. \textit{Spectral Analysis for Physical Applications: Multitaper and Conventional Univariate Techniques}. Cambridge University Press, 1993.

多窗谱估计与经典单变量谱分析技术的系统教材。书中对 Slepian 序列与 Thomson 多窗方法的完整处理，为本书时序分析章与离散采样逼近章的谱估计方法提供了标准参考。

\bibitem{oppenheim2010} \textbf{Alan V. Oppenheim, Ronald W. Schafer}. \textit{Discrete-Time Signal Processing}. 3rd ed., Pearson, 2010.

离散时间信号处理的标准教材，含周期卷积、循环卷积与 DFT 的详细处理。

\bibitem{mallat2009} \textbf{St\'ephane Mallat}. \textit{A Wavelet Tour of Signal Processing}. 3rd ed., Academic Press, 2009.

小波信号处理的权威著作。

\bibitem{daubechies1992} \textbf{Ingrid Daubechies}. \textit{Ten Lectures on Wavelets}. SIAM, 1992.

小波理论的经典教材。墨西哥帽核是小波变换的标准母函数。

\bibitem{cooley1965} \textbf{James W. Cooley, John W. Tukey}. \textit{An Algorithm for the Machine Calculation of Complex Fourier Series}. \textit{Mathematics of Computation}, 19(90), 297--301 (1965).

Cooley-Tukey快速傅里叶变换（FFT）的原始论文。

\bibitem{bluestein1970} \textbf{Leo I. Bluestein}. \textit{A Linear Filtering Approach to the Computation of the Discrete Fourier Transform}. \textit{IEEE Transactions on Audio and Electroacoustics}, 18(4), 451--455 (1970).

Bluestein算法的原始论文。

\bibitem{marks1991} \textbf{Robert J. Marks II}. \textit{Introduction to Shannon Sampling and Interpolation Theory}. Springer, 1991.

香农采样与插值理论的入门教材，含采样定理的推广与基数级数。

\bibitem{elder2015} \textbf{Yonina C. Eldar}. \textit{Sampling Theory: Beyond Bandlimited Systems}. Cambridge University Press, 2015.

超越带限系统的采样理论，含次奈奎斯特采样与压缩感知。

\bibitem{khinchin1963} \textbf{A. Ya. Khinchin}. \textit{Continued Fractions}. 3rd ed., P. Noordhoff, 1963.

连分数理论的经典著作。书中对收敛子的最佳逼近性质、逼近阶估计与丢番图逼近的度量理论的系统处理，构成本书离散采样逼近与置闰章中置闰误差有界定理与连分数最优性的数学基础。

\bibitem{lang1995} \textbf{Serge Lang}. \textit{Introduction to Diophantine Approximations}. Springer, 1995.

丢番图逼近的入门教材，含连分数、最佳逼近与克罗内克逼近定理。

\bibitem{walters1981} \textbf{Peter Walters}. \textit{An Introduction to Ergodic Theory}. Springer, 1981.

遍历理论的经典教材。书中对 Birkhoff 遍历定理、Poincar\'e 回复定理与紧致度量空间上连续变换的遍历理论的系统阐述，为本书概率论与分布章中遍历性与时间平均-空间平均判据提供了标准参考。

\bibitem{krengel1985} \textbf{Ulrich Krengel}. \textit{Ergodic Theorems}. de Gruyter, 1985.

遍历定理的权威专著，含 von Neumann 平均遍历定理与 Birkhoff 逐点遍历定理的系统证明。

\bibitem{lee2012smooth} \textbf{John M. Lee}. \textit{Introduction to Smooth Manifolds}. 2nd ed., Springer, 2012.

光滑流形理论的现代标准教材。书中对流形、切丛、向量场与微分形式的系统处理，为本书流形基础与纤维丛章、周期流形章提供了标准参考。

\bibitem{wang2023manifolds} \textbf{王正栋}. 《流形与纤维丛》。待补出版信息。

本书流形基础与纤维丛章的基础教材。

\bibitem{nakahara2003} \textbf{Mikio Nakahara}. \textit{Geometry, Topology and Physics}. 2nd ed., IOP Publishing, 2003.

规范场论的纤维丛表述。

\bibitem{conner1979} \textbf{Pierre E. Conner}. \textit{Differentiable Periodic Maps}. Springer, 1979.

微分周期映射的拓扑理论，为本书周期流形章提供周期性在流形上的拓扑背景。

\bibitem{atiyahsinger1968} \textbf{Michael Atiyah, Isadore M. Singer}. \textit{The Index of Elliptic Operators}. \textit{Annals of Mathematics}, 87(3), 484--530 (1968).

Atiyah-Singer指标定理。

\bibitem{bohm2003} \textbf{Arno Bohm, Ali Mostafazadeh, Hiroyasu Koizumi, Qian Niu, Joseph Zwanziger}. \textit{The Geometric Phase in Quantum Systems: Foundations, Mathematical Concepts, and Applications in Molecular and Condensed Matter Physics}. Springer, 2003.

几何相位的系统专著。书中对 Berry 绝热相位、几何联络与曲率、以及几何相位在分子与凝聚态物理中的应用的系统论述，为本书电磁系统章与周期流形章中相位的几何解释提供了标准参考。

\bibitem{feller1950} \textbf{William Feller}. \textit{An Introduction to Probability Theory and Its Applications}, Vol.~1. Wiley, 1950.

概率论的经典教材。书中对随机变量、分布函数、特征函数与极限定理的系统处理，为本书概率论与分布章提供了标准参考。

\bibitem{feller1966} \textbf{William Feller}. \textit{An Introduction to Probability Theory and Its Applications}, Vol.~2. Wiley, 1966.

概率论第二卷，含测度论、随机游走、拉普拉斯变换与概率分布的系统处理。

\bibitem{onsager1944} \textbf{Lars Onsager}. \textit{Crystal Statistics. I. A Two-Dimensional Model with an Order-Disorder Transition}. \textit{Physical Review}, 65(3--4), 117--149 (1944).

Onsager二维Ising模型的严格解，是相变与临界现象理论的奠基之作。

\bibitem{boltzmann1877} \textbf{Ludwig Boltzmann}. \textit{\"Uber die Beziehung zwischen dem zweiten Hauptsatze der mechanischen W\"armetheorie und der Wahrscheinlichkeitsrechnung}. \textit{Wiener Berichte}, 76, 373--435 (1877).

Boltzmann 原理 $S = k_B \ln W$ 的原始文献。该原理将熵与相空间中恒定能量流形的几何面积联系起来，为本书概率论与分布章中熵与分布的关系提供了历史与理论基础。

\bibitem{landau1987} \textbf{L. D. Landau, E. M. Lifshitz}. \textit{Fluid Mechanics}. 2nd ed., Course of Theoretical Physics, Vol.~6, Pergamon Press, 1987.

流体力学经典教材。书中对理想流体与粘性流体的运动方程、湍流与边界层的系统论述，为本书地理系统章与气候系统章中流体动力学背景提供了标准参考。

\bibitem{arnoldkhesin1998} \textbf{Vladimir I. Arnold, Boris A. Khesin}. \textit{Topological Methods in Hydrodynamics}. Springer, 1998.

拓扑流体力学领域的权威专著。书中从统一观点处理理想流体与磁流体力学的拓扑、群表示论与几何问题，为本书周期流形章与地理系统章中流体运动的拓扑约束提供了标准参考。

\bibitem{nocedal2006} \textbf{Jorge Nocedal, Stephen J. Wright}. \textit{Numerical Optimization}. 2nd ed., Springer, 2006.

数值优化的标准教材，含梯度下降、牛顿法与拟牛顿法的系统分析。

\bibitem{kochenderfer2019} \textbf{Mykel J. Kochenderfer, Tim A. Wheeler}. \textit{Algorithms for Optimization}. MIT Press, 2019.

优化算法的实用教材，含导数、梯度与多维度优化。

\bibitem{boyd2004} \textbf{Stephen P. Boyd, Lieven Vandenberghe}. \textit{Convex Optimization}. Cambridge University Press, 2004.

凸优化的标准教材。

\bibitem{nesterov2018} \textbf{Yurii Nesterov}. \textit{Lectures on Convex Optimization}. 2nd ed., Springer, 2018.

凸优化的经典讲义。Nesterov加速梯度法为KDI中多窗口加权合成的加速策略提供了理论参考。

\bibitem{goodfellow2016} \textbf{Ian Goodfellow, Yoshua Bengio, Aaron Courville}. \textit{Deep Learning}. MIT Press, 2016.

深度学习的标准教材，含反向传播与优化动力学的系统讨论。

\bibitem{hirose2006} \textbf{Akira Hirose}. \textit{Complex-Valued Neural Networks}. Springer, 2006.

复值神经网络的首部专著，为本书复指数激活函数（CEAF）提供神经网络语境下的数学基础。

\bibitem{hirose2012} \textbf{Akira Hirose} (ed.). \textit{Complex-Valued Neural Networks: Advances and Applications}. Wiley-IEEE Press, 2012.

复值神经网络的进展与应用，含幅相型激活函数与反向传播的复值推广。

\bibitem{scholkopf2002} \textbf{Bernhard Sch\"olkopf, Alexander J. Smola}. \textit{Learning with Kernels}. MIT Press, 2002.

核机器学习的权威著作。再生核希尔伯特空间理论为核函数变换矩阵提供了严格基础。

\bibitem{jackson1999} \textbf{John D. Jackson}. \textit{Classical Electrodynamics}. 3rd ed., Wiley, 1999.

经典电动力学的权威教材。书中关于电磁系统周期结构、麦克斯韦方程组导出以及洛伦兹力公式的内容为本研究提供了理论参考。

\bibitem{griffiths2013} \textbf{David J. Griffiths}. \textit{Introduction to Electrodynamics}. 4th ed., Pearson, 2013.

电动力学的标准入门教材。

\bibitem{alexander2017} \textbf{Charles K. Alexander, Matthew N. O. Sadiku}. \textit{Fundamentals of Electric Circuits}. 6th ed., McGraw-Hill, 2017.

电路分析的标准教材。

\bibitem{razavi2017} \textbf{Behzad Razavi}. \textit{Design of Analog CMOS Integrated Circuits}. 2nd ed., McGraw-Hill, 2017.

模拟CMOS集成电路设计的权威教材，涵盖晶振、PLL和时钟树的设计理论。

\bibitem{weste2011} \textbf{Neil H. E. Weste, David Money Harris}. \textit{CMOS VLSI Design: A Circuits and Systems Perspective}. 4th ed., Addison-Wesley, 2011.

VLSI设计的标准教材。

\bibitem{darwin1859} \textbf{Charles Darwin}. \textit{On the Origin of Species}. John Murray, 1859.

自然选择理论的原创著作。

\bibitem{lotka1925} \textbf{Alfred J. Lotka}. \textit{Elements of Physical Biology}. Williams and Wilkins, 1925.

Lotka-Volterra捕食者-猎物方程的原始出处。

\bibitem{volterra1926} \textbf{Vito Volterra}. \textit{Variazioni e fluttuazioni del numero d'individui in specie animali conviventi}. \textit{Memorie della Regia Accademia Nazionale dei Lincei}, 2, 31--113 (1926).

Volterra独立于Lotka推导出捕食者-猎物方程。

\bibitem{winfree2001} \textbf{Arthur T. Winfree}. \textit{The Geometry of Biological Time}. 2nd ed., Springer, 2001.

生物时间的几何理论，处理节律性返回循环的数学结构。



非线性动力学与混沌的经典教材，极限环与周期振荡的标准参考。

\bibitem{aschoff1965} \textbf{J\"urgen Aschoff}. \textit{Circadian Rhythms in Man}. \textit{Science}, 148(3676), 1427--1432 (1965).

人体昼夜节律的经典研究。

\bibitem{forger2017} \textbf{Daniel B. Forger}. \textit{Biological Clocks, Rhythms, and Oscillations: The Theory of Biological Timekeeping}. MIT Press, 2017.

生物钟、节律与振荡的数学建模教材，含昼夜节律的极限环分析与相位响应曲线。

\bibitem{nowak2006} \textbf{Martin A. Nowak}. \textit{Evolutionary Dynamics}. Harvard University Press, 2006.

进化动力学的数学理论。

\bibitem{kondratiev1935} \textbf{Nikolai D. Kondratiev}. \textit{The Long Waves in Economic Life}. \textit{Review of Economic Statistics}, 17(6), 105--115 (1935).

康德拉季耶夫长波理论。

\bibitem{kennedy1987} \textbf{Paul Kennedy}. \textit{The Rise and Fall of the Great Powers}. Random House, 1987.

大国兴衰的经典著作。

\bibitem{toynbee1934} \textbf{Arnold J. Toynbee}. \textit{A Study of History} (12 vols). Oxford University Press, 1934--1961.

汤因比的挑战-应战理论。

\bibitem{braudel1949} \textbf{Fernand Braudel}. \textit{La M\'editerran\'ee et le Monde M\'editerran\'een}. Armand Colin, 1949.

布罗代尔的多时间尺度历史观。

\bibitem{maclane1998} \textbf{Saunders Mac Lane}. \textit{Categories for the Working Mathematician}. 2nd ed., Springer, 1998.

范畴论的标准教材，本书周期范畴章的函子与自然变换的理论基础。

\bibitem{awodey2010} \textbf{Steve Awodey}. \textit{Category Theory}. 2nd ed., Oxford University Press, 2010.

范畴论的现代入门教材。

\bibitem{ren2011} \textbf{Wei Ren, Yongcan Cao}. \textit{Distributed Coordination of Multi-agent Networks}. Springer, 2011.

多智能体网络的分布式协调，含集体周期运动协调与周期解。

\bibitem{mackay2003} \textbf{David J. C. MacKay}. \textit{Information Theory, Inference, and Learning Algorithms}. Cambridge University Press, 2003.

信息论与贝叶斯推断的交叉领域杰作。

\bibitem{cover2006} \textbf{Thomas M. Cover, Joy A. Thomas}. \textit{Elements of Information Theory}. 2nd ed., Wiley, 2006.

信息论的标准教材。

\bibitem{jeffreys1961} \textbf{Harold Jeffreys}. \textit{Theory of Probability}. 3rd ed., Oxford University Press, 1961.

贝叶斯统计的经典著作。

\bibitem{sutton2018} \textbf{Richard S. Sutton, Andrew G. Barto}. \textit{Reinforcement Learning: An Introduction}. 2nd ed., MIT Press, 2018.

强化学习中奖励信号与策略优化的框架。

\bibitem{levinpereswilmer2017} \textbf{David A. Levin, Yuval Peres, Elizabeth L. Wilmer}. \textit{Markov Chains and Mixing Times}. 2nd ed., American Mathematical Society, 2017.

马尔可夫链混合时间的标准教材；第 12 章给出谱隙、$\chi^2$ 距离收缩与总变差距离界（定理 12.3、推论 12.5）。

\bibitem{kim2026} \textbf{Eloy Kim}. \textit{Periodicity Mathematics: A Unified Framework of Epicycle Series and Interdisciplinary Periodic Structures}. Manuscript, 2026.

本书的完整手稿。建立了周期数学的统一框架，包含周期基础、数系、标量与矢量、函数与分析基础、级数与卷积、微分方程与复变函数、线性代数、流形基础与纤维丛、周期流形、周期范畴、概率论、矢量分析、群与群函数、高阶群与高阶函数、周期函数与具有周期轨道的群、调和分析、时序分析、周期系统等核心章节。

\end{thebibliography}

\endgroup


\vspace{1cm}
\begin{flushright}
\textit{—— 谨以本书献给所有追寻“道”与“数”的探索者。}
\end{flushright}

\cleardoublepage
\phantomsection
\addcontentsline{toc}{chapter}{作者简介}
\markboth{作者简介}{作者简介}
\chapter*{作者简介}

Eloy Kim，独立研究者，长期从事周期数学、调和分析与跨学科数学建模研究。作者熟练使用 C/C++、Java、TypeScript/JavaScript 及 Android 开发，兼具工程实现与理论推导能力。研究兴趣涵盖《易经》、编程、数学与物理学，致力于在哲学与数理之间寻找统一的结构与语言。

\vspace{1em}
\noindent\textbf{联系方式：} \texttt{eloy.kim@outlook.com}

\cleardoublepage
\phantomsection
\addcontentsline{toc}{chapter}{致谢}
\markboth{致谢}{致谢}
\chapter*{致谢}

本研究的顺利完成，离不开众多师长、友人与同行所提供的宝贵研究资金支持与物质帮助。谨此致以诚挚的谢意。他们的资助为本书所涉的文献整理、数学推导与跨学科案例研究提供了必要的物质保障，使作者能够集中精力完成这一长期的研究纲领。

下列表格列出了已知姓名或机构的主要资助者与支持者（按致谢顺序，不分先后）。作者将随着本纲领的推进持续补充此表。

\begin{center}
\renewcommand{\arraystretch}{1.35}
\setlength{\tabcolsep}{4pt}
\small
\begin{longtable}{@{}>{\raggedright\arraybackslash}p{1.1cm}
                     >{\raggedright\arraybackslash}p{3.0cm}
                     >{\raggedright\arraybackslash}p{3.4cm}
                     >{\raggedright\arraybackslash}p{5.6cm}@{}}
\caption*{}\\
\toprule
\textbf{序号} & \textbf{资助者/支持者} & \textbf{支持类型} & \textbf{致谢备注} \\
\midrule
\endfirsthead

\toprule
\textbf{序号} & \textbf{资助者/支持者} & \textbf{支持类型} & \textbf{致谢备注} \\
\midrule
\endhead

\midrule
\multicolumn{4}{r}{\small 接下页} \\
\endfoot

\bottomrule
\endlastfoot

1 & HYC & 物质与精神支持 & 在本书写作的多年间持续提供物质支持与精神鼓励，是研究得以坚持的重要依托。 \\
2 & Tong Su & 研究资金支持 & 在研究纲领的起始阶段给予关键资助与思想启发，使本研究得以启动。后续虽联系渐少，但最初的推动不可替代。 \\
3 & JD Li & 研究资金支持 & 持续资助跨学科数学探索，为文献整理与案例研究提供了稳定的物质保障。 \\
4 & 匿名支持者 & 物质支持 & 在本书早期阶段给予物质支持。作者因故未能获知其姓名，谨此致谢。 \\
5 & 中途离开的支持者 & 物质支持 & 曾在研究纲领早期给予资金或物质帮助，后因故中断联系。作者铭记其贡献，谨致谢意。 \\

\end{longtable}
\end{center}

此外，还有一批在本书写作的不同阶段给予过资金、资料、设备、场地或其他物质支持、但作者未能一一记下其姓名或不愿公开姓名的朋友与机构。他们的帮助同样不可或缺，谨此一并致以诚挚的谢意。

感谢所有对基础研究与跨学科探索给予信任与支持的人。

\vspace{1em}
\begin{flushright}
作者Eloy Kim\\
\today
\end{flushright}

\vspace{2cm}
\begin{center}
    \textbf{周期数学}\\
    \copyright\ 2026 \\
    Licence: Apache 2.0
\end{center}

